% ══════════════════════════════════════════════════════════════════════
% Viewpoint 4 Corrigendum: Cercis-Fisher is Conditional Equivalence, NOT Natural Emergence
% 观点4更正：Cercis-Fisher是条件等价，非自然涌现
%
% This is a standalone corrigendum to be inserted after Section 4 (Information-Geometric
% Bulk-Boundary Correspondence) of gauge_formalized.tex, or read alongside it.
%
% Based on formal peer review (docs/reviews/GAUGE_5VIEWPOINTS_FINAL.md,
% docs/reviews/GAUGE_VIEWPOINTS_REVIEW.md).
% ══════════════════════════════════════════════════════════════════════

\documentclass{article}
\usepackage[UTF8]{ctex}
\usepackage{amsmath,amssymb,amsthm}
\usepackage{hyperref}
\usepackage{enumitem}
\usepackage[margin=2.5cm]{geometry}

% ─── Math notation (consistent with gauge_formalized.tex) ───
\newcommand{\GeoDist}{\mathrm{GeoDist}}
\newcommand{\KL}{\mathrm{KL}}
\newcommand{\Fisher}{\mathcal{I}}
\newcommand{\cercis}{\mathcal{C}}
\newcommand{\cercisnorm}{\bar{\cercis}}
\newcommand{\Situs}{\mathcal{S}}
\newcommand{\ExpFam}{\mathcal{E}}
\newcommand{\Od}{\mathrm{O}(d)}
\newcommand{\SOthree}{\mathrm{SO}(3)}
\newcommand{\ZBetti}{\beta_1}
\newcommand{\Lone}{L_1}
\newcommand{\edgs}{\mathcal{E}}
\newcommand{\verts}{\mathcal{V}}
\newcommand{\face}{\mathcal{F}}
\newcommand{\grph}{\mathcal{G}}
\newcommand{\Gflatquot}{\mathcal{M}_{\mathrm{flat}}}
\newcommand{\triv}{\mathbf{1}}
\newcommand{\doneT}{d_1^{\mathsf{T}}}
\DeclareMathOperator{\im}{im}
\DeclareMathOperator{\Log}{Log}
\DeclareMathOperator{\Ad}{Ad}
\DeclareMathOperator{\Hom}{Hom}
\DeclareMathOperator{\Stab}{Stab}
\newcommand{\E}{\mathbb{E}}
\DeclareMathOperator{\dist}{dist}

\title{\textbf{观点4更正（Corrigendum）} \\
\large Cercis-Fisher等价性：条件等价而非自然涌现 \\
\large Cercis-Fisher Equivalence: Conditional Equivalence, Not Natural Emergence}
\author{SCX 理论架构组 / SCX Theory Architecture Group}
\date{\today}

\begin{document}
\maketitle

% ══════════════════════════════════════════════════════════════════════
\section{审查裁决 / Review Verdict}
% ══════════════════════════════════════════════════════════════════════

\begin{center}
\begin{tabular}{|p{6cm}|p{8cm}|}
\hline
\textbf{原主张 / Original Claim} & \textbf{裁决 / Verdict} \\
\hline
Cercis${}^2 \propto$ 平均KL散度 $=$ Fisher测地线长度平方。
从信息几何自然涌现。
& ⚠️ 数学关系正确但``自然涌现''夸大 / Mathematically correct but ``natural emergence'' is overstated \\
\hline
\end{tabular}
\end{center}

\subsection{具体问题 / Specific Issues}

\begin{enumerate}[label=(\arabic*)]
    \item \textbf{不是``平均KL''而是``最小KL''。}
    Cercis是到纯规范子流形的最小KL散度：
    $\cercis \approx 2 \cdot \min_{h} \KL(P_A \| P_{d_0 h})$，
    而非``平均KL散度''。

    \textbf{Not ``average KL'' but ``minimum KL''.}
    Cercis is the \emph{minimum} KL divergence to the pure-gauge submanifold,
    not an average.

    \item \textbf{不是``自然涌现''而是条件等价。}
    $\GeoDist^2 \leftrightarrow 2\KL$ 的等价性依赖8个强假设（见第\ref{sec:assumptions}节）。
    任何一个假设不满足即断裂。SCX数据极不可能恰好形成指数族。

    \textbf{Not ``natural emergence'' but conditional equivalence.}
    The $\GeoDist^2 \leftrightarrow 2\KL$ equivalence depends on 8 strong assumptions (see \S\ref{sec:assumptions}).
    If any single assumption fails, the equivalence breaks. SCX data is extremely unlikely to form an exponential family.

    \item \textbf{非指数族下完全失效。}
    对一般统计流形，$\GeoDist^2$与$2\KL$的差由Amari-Chentsov三阶张量控制。
    该张量在一般SCX数据上非零，等价性由此退化。

    \textbf{Complete breakdown outside exponential families.}
    For general statistical manifolds, the difference between $\GeoDist^2$ and $2\KL$ is controlled
    by the Amari-Chentsov third-order tensor. This tensor is generically non-zero for SCX data,
    so the equivalence degrades.
\end{enumerate}

% ══════════════════════════════════════════════════════════════════════
\section{八大假设清单 / The Eight Assumptions}
\label{sec:assumptions}
% ══════════════════════════════════════════════════════════════════════

要使得 $\GeoDist(p,q)^2 = 2\KL(p\|q) + O(\|p-q\|^3)$ 严格成立，
以下八个假设必须同时满足。任何一个的失效都会导致等式退化。

For $\GeoDist(p,q)^2 = 2\KL(p\|q) + O(\|p-q\|^3)$ to hold rigorously,
the following eight assumptions must be simultaneously satisfied.
Failure of any single one degrades the equality.

\vspace{1em}

\begin{enumerate}[label=\textbf{H\arabic*:}, leftmargin=*, itemsep=0.8em]

    \item \textbf{指数族假设 (Exponential Family Assumption).} \\
    \textbf{中文：} 输出分布族 $\mathcal{P} = \{p(x \mid \theta)\}$ 必须构成指数族，
    即存在充分统计量 $T(x)$ 和累积量生成函数 $\psi(\theta)$ 使得
    $p(x \mid \theta) = h(x) \exp(\langle\theta, T(x)\rangle - \psi(\theta))$。
    只有在此条件下，Fisher信息矩阵才是 $\psi$ 的Hessian，
    KL散度才等于Bregman散度（引理4.3 / Lemma~\ref{lem:kl_bregman}的原上下文）。
    \\[0.3em]
    \textbf{English:} The output distribution family $\mathcal{P}$ must be an exponential family.
    Only then is the Fisher matrix the Hessian of $\psi$, and KL divergence equals Bregman divergence.
    \\[0.3em]
    \textbf{断裂后果：} 非指数族下，$\KL(\theta_1\|\theta_2)$ 不再是 $\psi$ 的Bregman散度，
    Taylor展开的第二项不再是 $g_{ij}\Delta\theta^i\Delta\theta^j$。
    $\GeoDist^2$ 与 $2\KL$ 的差由 Amari-Chentsov 三阶张量控制（开放问题6.2 / Open Problem~\ref{prob:non_exp_fam}）。
    \\[0.3em]
    \textbf{Consequence of failure:} Outside exponential families, KL is no longer a Bregman divergence,
    and the second-order Taylor term is no longer $g_{ij}\Delta\theta^i\Delta\theta^j$.
    The difference is controlled by the Amari-Chentsov third-order tensor (Open Problem~\ref{prob:non_exp_fam}).

    \item \textbf{小距离／局部近似假设 (Small-Distance / Local Approximation Assumption).} \\
    \textbf{中文：} $\theta_p$ 与 $\theta_q$ 必须充分接近，使得三阶及以上余项
    $O(\|\theta_q - \theta_p\|^3)$ 可忽略。该等式的所有推导均基于Taylor展开截断至二阶。
    \\[0.3em]
    \textbf{English:} $\theta_p$ and $\theta_q$ must be sufficiently close so that third-order
    and higher remainders $O(\|\theta_q - \theta_p\|^3)$ are negligible. All derivations
    truncate the Taylor expansion at second order.
    \\[0.3em]
    \textbf{断裂后果：} 在大分歧（large discrepancy）情形下，三阶误差项主导，
    $\GeoDist^2 \approx 2\KL$ 不再成立。SCX跨领域专家的大分歧是常见现象。
    \\[0.3em]
    \textbf{Consequence of failure:} For large discrepancies (common in cross-domain SCX experts),
    third-order error terms dominate and $\GeoDist^2 \approx 2\KL$ no longer holds.

    \item \textbf{$C^2$ 可微性假设 ($C^2$ Differentiability Assumption).} \\
    \textbf{中文：} 映射 $\theta \mapsto p(x \mid \theta)$ 必须对几乎所有 $x$ 是 $C^2$ 的。
    这是Taylor展开至二阶、以及Fisher信息矩阵作为二阶导数的定义的前提。
    \\[0.3em]
    \textbf{English:} The map $\theta \mapsto p(x \mid \theta)$ must be $C^2$ for almost all $x$.
    This is required for second-order Taylor expansion and for the Fisher matrix as a second derivative.
    \\[0.3em]
    \textbf{断裂后果：} 如果密度函数包含不可微点（如ReLU激活产生的分段分布），Fisher度量在这些点不定义，
    测地线无法唯一确定。
    \\[0.3em]
    \textbf{Consequence of failure:} If the density has non-differentiable points (e.g., piecewise
    distributions from ReLU activations), the Fisher metric is undefined at those points and
    geodesics are not uniquely determined.

    \item \textbf{积分-微分可交换假设 (Integration-Differentiation Interchangeability Assumption).} \\
    \textbf{中文：} $\partial_\theta \int p(x|\theta) d\mu(x) = \int \partial_\theta p(x|\theta) d\mu(x)$
    必须成立。这是Fisher信息矩阵的标准正则性条件，确保 $\E[\partial_\theta \log p] = 0$。
    \\[0.3em]
    \textbf{English:} The interchange $\partial_\theta \int p = \int \partial_\theta p$ must hold.
    This is the standard regularity condition ensuring $\E[\partial_\theta \log p] = 0$, which
    underpins the Fisher information definition.
    \\[0.3em]
    \textbf{断裂后果：} 若支撑集依赖于参数（如均匀分布），Fisher信息的标准形式不再成立。
    \\[0.3em]
    \textbf{Consequence of failure:} If the support depends on $\theta$ (e.g., uniform distributions),
    the standard Fisher information form breaks down.

    \item \textbf{Fisher满秩假设 (Fisher Full-Rank Assumption).} \\
    \textbf{中文：} Fisher信息矩阵 $\Fisher(\theta)$ 必须在每一点正定（满秩）。
    这是 $\Theta$ 上Fisher度量成为真正黎曼度量的必要条件——退化度量不定义唯一的测地线。
    \\[0.3em]
    \textbf{English:} The Fisher information matrix $\Fisher(\theta)$ must be positive-definite
    (full rank) at every point. This is required for $(\Theta, g)$ to be a genuine Riemannian
    manifold — a degenerate metric does not define unique geodesics.
    \\[0.3em]
    \textbf{断裂后果：} 若参数间存在函数依赖（如归一化约束导致的有效降维），Fisher矩阵奇异，
    测地线不存在或不唯一，$\GeoDist$ 无定义。
    \\[0.3em]
    \textbf{Consequence of failure:} If parameters are functionally dependent (e.g., normalization
    constraints causing effective dimensionality reduction), the Fisher matrix is singular,
    geodesics may not exist or be unique, and $\GeoDist$ is undefined.

    \item \textbf{小联络近似假设 (Small-Connection Approximation Assumption).} \\
    \textbf{中文：} 边赋值 $a_e$ 与规范变换 $h_v$ 必须在小联络极限内
    （$\max_e \dist_{\Od}(A_e, I_d) < r_0$，其中 $r_0$ 为矩阵对数的收敛半径）。
    只有在此极限下，非阿贝尔规范固定才退化为李代数 $\mathfrak{so}(d)$ 上的线性霍奇问题。
    \\[0.3em]
    \textbf{English:} Edge assignments $a_e$ and gauge transformations $h_v$ must be within the
    small-connection limit ($\max_e \dist_{\Od}(A_e, I_d) < r_0$, the matrix logarithm convergence radius).
    Only then does non-abelian gauge-fixing reduce to linear Hodge theory on $\mathfrak{so}(d)$.
    \\[0.3em]
    \textbf{断裂后果：} 大旋转下（如跨模态专家间的$90^\circ$标架差异），Cartan分解非线性项不能忽略，
    $\GeoDist$ 与欧氏距离的关系断裂。
    \\[0.3em]
    \textbf{Consequence of failure:} Under large rotations (e.g., $90^\circ$ frame misalignment
    between cross-modal experts), nonlinear Cartan terms cannot be neglected, and the relationship
    between $\GeoDist$ and Euclidean distance breaks.

    \item \textbf{欧氏近似假设 (Euclidean Approximation / Degenerate Fisher Assumption).} \\
    \textbf{中文：} Cercis的欧氏定义（$\|r\|^2 = \|r_{\text{harm}}\|^2 + \|r_{\text{coexact}}\|^2$）
    要与Fisher定义等价，需要Fisher度量退化为欧氏度量：$\Fisher_{ij}(\theta) = \delta_{ij}$。
    即信息流形的所有方向权重相等，曲率为零。
    \\[0.3em]
    \textbf{English:} For the Euclidean Cercis definition ($\|r\|^2 = \|r_{\text{harm}}\|^2 + \|r_{\text{coexact}}\|^2$)
    to coincide with the Fisher definition, the Fisher metric must degenerate to Euclidean:
    $\Fisher_{ij}(\theta) = \delta_{ij}$. This means all directions on the information manifold
    have equal weight and zero curvature.
    \\[0.3em]
    \textbf{断裂后果：} 真实SCX数据的参数空间几乎肯定具有非平凡的Fisher曲率。
    Cercis的欧氏范数和Fisher测地距离之间存在不可忽略的曲率修正。
    \\[0.3em]
    \textbf{Consequence of failure:} Real SCX parameter spaces almost certainly have non-trivial
    Fisher curvature. A non-negligible curvature correction exists between Euclidean Cercis and
    Fisher geodesic distance.

    \item \textbf{分布紧致假设 (Distribution Compactness / Bounded Support Assumption).} \\
    \textbf{中文：} 所有专家输出分布必须在 $\Theta$ 的紧致子集上充分集中，使得局部Taylor展开
    在边界观测和所有规范体态上一致有效。
    \\[0.3em]
    \textbf{English:} All expert output distributions must be sufficiently concentrated on a compact
    subset of $\Theta$, so that the local Taylor expansion is uniformly valid across all
    boundary observations and gauge bulk states.
    \\[0.3em]
    \textbf{断裂后果：} 重尾分布（如学生-t分布）可能使某些观测远离局部有效域，
    导致三阶误差项在个别边上不可控地增长。
    \\[0.3em]
    \textbf{Consequence of failure:} Heavy-tailed distributions (e.g., Student's-t) may place
    some observations outside the local validity domain, causing uncontrolled growth of
    third-order error terms on individual edges.
\end{enumerate}

% ══════════════════════════════════════════════════════════════════════
\section{修正后的主张 / Corrected Claim}
% ══════════════════════════════════════════════════════════════════════

\subsection{中文修正 (Chinese Correction)}

{\bf 原主张（夸大，需删除）：}
\begin{quote}
``Cercis分数是Fisher信息度量下的测地线距离。指数族假设下Cercis${}^2 \propto$ 平均KL散度
$=$ Fisher测地线长度平方。审计分数不是人为定义——从数据流形信息几何自然涌现。''
\end{quote}

{\bf 修正后的主张（条件等价）：}
\begin{quote}
{\bf 在指数族假设和局部近似下，}Cercis分数等于观测边分布到纯规范子流形的{\bf 最小}KL散度的两倍，
而后者又等于（至二阶）Fisher测地线距离的平方。
这提供了一种信息几何{\bf 解释}——但并非``自然涌现''，而是{\bf 条件等价}：
当且仅当Situs流形承认指数族结构时成立。
{\bf 对于非指数族数据，该关系由Amari-Chentsov三阶张量退化（开放问题6.2）。}
\end{quote}

具体表述为：

\vspace{0.5em}
\fbox{%
\begin{minipage}{0.94\textwidth}
\vspace{0.8em}
\noindent {\bf 定理（Cercis的信息几何条件等价性，修订版）}
\label{thm:cercis_conditional_equivalence}

{\bf 假设：} 上述8个假设(H1--H8)同时成立。

{\bf 结论：}
\[
\cercis = \min_{d_0 h \in \im(d_0)} \GeoDist(P_A,\; P_{d_0 h})^2
       = 2 \min_{d_0 h \in \im(d_0)} \KL(P_A \| P_{d_0 h}) + O(\|\mathfrak{a}\|^3).
\]

Cercis度量了从观测专家边缘分布到``纯规范体态''（可完美对齐的假想世界）
的最小Fisher信息距离。

{\bf 限制：}
\begin{enumerate}[label=(\roman*)]
    \item 该等价性在指数族(H1)内严格成立；非指数族下，$\GeoDist^2$与$2\KL$的差
    由Amari-Chentsov三阶张量$T_{ijk} = \E[\partial_i \log p \cdot \partial_j \log p \cdot \partial_k \log p]$
    控制，且$T_{ijk}$在一般SCX数据上非零（开放问题\ref{prob:non_exp_fam}）。
    \item 该等价性在局部(H2)近似下成立；大分歧下三阶项主导。
    \item ``自然涌现''是错误的描述——这是在一组强假设下的构造性等价，
    而非数据固有的几何必然性。
\end{enumerate}
\vspace{0.5em}
\end{minipage}%
}
\vspace{0.5em}

\vspace{1em}

\subsection{English Correction}

{\bf Original claim (overstated, to be removed):}
\begin{quote}
``The Cercis Score is the geodesic distance under the Fisher information metric.
Under the exponential family assumption, Cercis${}^2 \propto$ mean KL divergence
$=$ squared Fisher geodesic length. The audit score is not artificially defined —
it naturally emerges from the information geometry of the data manifold.''
\end{quote}

{\bf Corrected claim (conditional equivalence):}
\begin{quote}
{\bf Under the exponential family assumption and local approximation,} the Cercis Score
equals twice the {\bf minimum} KL divergence from the observed edge distribution to the
pure-gauge submanifold, which in turn equals (to second order) the squared Fisher geodesic
distance. This provides an information-geometric {\bf interpretation} — not a ``natural emergence''
but a {\bf conditional equivalence} that holds when and only when the Situs manifold
admits an exponential family structure.
{\bf For non-exponential-family data, the relationship degrades by the Amari-Chentsov
tensor (Open Problem~6.2).}
\end{quote}

Explicit formulation:

\vspace{0.5em}
\fbox{%
\begin{minipage}{0.94\textwidth}
\vspace{0.8em}
\noindent {\bf Theorem (Information-Geometric Conditional Equivalence of Cercis, Revised)}
\label{thm:cercis_conditional_equivalence_en}

{\bf Assumptions:} All 8 assumptions (H1--H8) hold simultaneously.

{\bf Conclusion:}
\[
\cercis = \min_{d_0 h \in \im(d_0)} \GeoDist(P_A,\; P_{d_0 h})^2
       = 2 \min_{d_0 h \in \im(d_0)} \KL(P_A \| P_{d_0 h}) + O(\|\mathfrak{a}\|^3).
\]

Cercis measures the minimal Fisher information distance from the observed expert edge
distributions to the ``pure-gauge bulk states'' — the hypothetical world where all
experts can be perfectly aligned.

{\bf Limitations:}
\begin{enumerate}[label=(\roman*)]
    \item The equivalence is rigorous within exponential families (H1); outside them,
    the difference between $\GeoDist^2$ and $2\KL$ is controlled by the Amari-Chentsov
    third-order tensor $T_{ijk} = \E[\partial_i \log p \cdot \partial_j \log p \cdot \partial_k \log p]$,
    which is generically non-zero for SCX data (Open Problem~\ref{prob:non_exp_fam}).
    \item The equivalence holds under local approximation (H2); for large discrepancies,
    third-order terms dominate.
    \item ``Natural emergence'' is an incorrect characterization — this is a
    \emph{constructed equivalence} under strong assumptions, not an intrinsic
    geometric inevitability of the data.
\end{enumerate}
\vspace{0.5em}
\end{minipage}%
}
\vspace{0.5em}

% ══════════════════════════════════════════════════════════════════════
\section{五大断裂点 / Five Breakdown Points}
% ══════════════════════════════════════════════════════════════════════

以下总结了等价性在实际SCX数据上的五个主要断裂途径：

Below, we summarize five principal ways the equivalence can break on real SCX data:

\vspace{0.5em}

\begin{enumerate}[label=\textbf{B\arabic*:}, leftmargin=*, itemsep=0.8em]

    \item \textbf{非指数族失效（最严重） / Non-Exponential-Family Failure (Most Severe).} \\
    SCX专家输出（特别是MoE路由权重、softmax得分、注意力分布）极不可能恰好构成指数族。
    指数族要求特定的代数结构（线性充分统计量$T(x)$），而深度学习输出通常不满足。
    文档自身将其列为开放问题6.2（\ref{prob:non_exp_fam}）。
    \\[0.3em]
    SCX expert outputs (MoE routing weights, softmax scores, attention distributions)
    are extremely unlikely to form an exponential family. Exponential families require specific
    algebraic structure (linear sufficient statistics $T(x)$) that deep learning outputs do not
    generally satisfy. The document itself lists this as Open Problem~6.2 (\ref{prob:non_exp_fam}).

    \item \textbf{混淆``最小KL''与``平均KL'' / Confusing ``min KL'' with ``average KL''.} \\
    Cercis是最小化问题 $\min_h \KL(P_A \| P_{d_0 h})$ 的解，不是对所有专家对取平均的KL散度。
    这两个概念在数学上有本质区别。
    \\[0.3em]
    Cercis solves $\min_h \KL(P_A \| P_{d_0 h})$ — a minimization problem, not an average of
    KL divergences over all expert pairs. These are mathematically distinct.

    \item \textbf{欧氏近似的隐性要求 / Hidden Euclidean Approximation Requirement.} \\
    Cercis的欧氏范数定义与Fisher测地距离的等价性要求 $\Fisher_{ij} = \delta_{ij}$ (H7)。
    这意味着信息流形必须处处平坦——在实际的SCX参数空间中极不可能成立。
    \\[0.3em]
    The equivalence between Euclidean-norm Cercis and Fisher geodesic distance requires
    $\Fisher_{ij} = \delta_{ij}$ (H7) — meaning the information manifold must be everywhere flat,
    which is extremely unlikely on real SCX parameter spaces.

    \item \textbf{定理4.2证明中的内在修正 / Intrinsic Correction in Theorem~4.2 Proof.} \\
    选错子流形会丢失调和分量。正确定义为：体态子流形由正合联络（纯规范）$\im(d_0)$ 参数化，
    而非调和体态 $\ker(\Lone)$。原定理\ref{thm:cercis_bulk_boundary}在证明中发现并修正了此问题。
    \\[0.3em]
    Choosing the wrong submanifold loses the harmonic component. The correct definition is:
    the bulk submanifold is parameterized by exact connections (pure gauge) $\im(d_0)$,
    not harmonic bulk states $\ker(\Lone)$. Theorem~\ref{thm:cercis_bulk_boundary} identifies
    and corrects this issue in its proof.

    \item \textbf{局部性限制 / Locality Constraint.} \\
    所有结果仅在 $\|\theta_q - \theta_p\|$ 充分小时成立。当跨构型的专家差异较大
    （如不同领域、不同训练阶段的专家），三阶误差项 $O(\|\Delta\theta\|^3)$ 不可忽略。
    \\[0.3em]
    All results hold only when $\|\theta_q - \theta_p\|$ is sufficiently small.
    When inter-configuration expert differences are large (e.g., different domains, different
    training stages), the third-order error term $O(\|\Delta\theta\|^3)$ is non-negligible.

\end{enumerate}

% ══════════════════════════════════════════════════════════════════════
\section{修订建议 / Recommended Revisions}
% ══════════════════════════════════════════════════════════════════════

\subsection{在主文件中需要的修改 / Changes Needed in gauge\_formalized.tex}

\begin{enumerate}[label=(\arabic*)]

    \item \textbf{摘要（第234--238行）：}
    删除``自然涌现''相关措辞。将
    \begin{quote}
        ``我们证明在此族下，Fisher测地线长度与KL散度满足局部等价性...''
    \end{quote}
    改为：
    \begin{quote}
        ``我们证明在指数族假设和局部近似下，Fisher测地线长度与KL散度之间存在条件等价关系...''
    \end{quote}
    英文摘要同理修改。

    \item \textbf{第\ref{sec:info}节标题：}
    将
    \begin{quote}
        ``指数族下的Fisher测地线与KL散度的等价性 / Fisher Geodesic--KL Equivalence under Exponential Families''
    \end{quote}
    改为：
    \begin{quote}
        ``指数族下的Fisher测地线与KL散度的条件等价性 / Conditional Fisher Geodesic--KL Equivalence under Exponential Families''
    \end{quote}

    \item \textbf{定理\ref{thm:cercis_bulk_boundary}：}
    在定理陈述之前添加明确的假设列表，引用上述H1--H8。在定理结尾添加限制性注记。

    \item \textbf{定理\ref{thm:unified}（统一结构定理）：}
    将措辞从
    \begin{quote}
        ``Cercis重新表述为从边界到纯规范体态的Fisher测地距离，在指数族下局部等价于$2\KL$''
    \end{quote}
    改为：
    \begin{quote}
        ``在指数族假设和局部近似下，Cercis可重新解释为从边界到纯规范体态的Fisher测地距离，
        条件等价于$2\KL$（至二阶）。非指数族下的行为由开放问题6.2处理。''
    \end{quote}

    \item \textbf{比较表（第1529--1532行）：}
    将
    \begin{quote}
        ``Cercis定义: $\min_{d_0 h} \GeoDist(P_A, P_{d_0 h})^2$'' 和
        ``Cercis概率解释: $\approx 2\min_{d_0 h} \KL(P_A \| P_{d_0 h})$''
    \end{quote}
    添加脚注或括号注：``（需指数族假设，非指数族下退化）''

    \item \textbf{第\ref{sec:comparison}节``信息几何的严格性''段落（第1560--1564行）：}
    将
    \begin{quote}
        ``本文将其应用于SCX的$\Situs$流形是自然的推广...''
    \end{quote}
    改为：
    \begin{quote}
        ``本文将其应用于SCX的$\Situs$流形是在特定假设下的条件推广。
        当这些假设不满足时（特别是非指数族情形，开放问题6.2），等价性由Amari-Chentsov张量退化。''
    \end{quote}

    \item \textbf{结论第3点（第1678--1684行）：}
    将
    \begin{quote}
        ``证明了指数族下Fisher测地线长度与KL散度的局部等价性...建立了Cercis分数的信息几何解释''
    \end{quote}
    改为：
    \begin{quote}
        ``证明了指数族假设下Fisher测地线长度与KL散度的条件等价性...建立了Cercis分数的条件信息几何解释，并明确标注了该解释的8个前提假设和失效条件。''
    \end{quote}

\end{enumerate}

% ══════════════════════════════════════════════════════════════════════
\section{开放问题6.2的交叉引用 / Cross-Reference to Open Problem 6.2}
% ══════════════════════════════════════════════════════════════════════

本更正与开放问题\ref{prob:non_exp_fam}（``非指数族下的Fisher-KL等价性''）直接关联。
该问题指出：

\begin{quote}
对于一般的SCX输出分布（可能非指数族），Fisher测地距离和KL散度之间的关系由
Amari-Chentsov三阶张量 $T_{ijk} = \E[\partial_i \log p \cdot \partial_j \log p \cdot \partial_k \log p]$
控制。需要该张量在SCX输出的经验分布上的界。
\end{quote}

This corrigendum is directly linked to Open Problem~\ref{prob:non_exp_fam}
(``Fisher-KL Equivalence Outside Exponential Families''), which states:

\begin{quote}
For general SCX output distributions (potentially non-exponential-family),
the relationship between Fisher geodesic distance and KL divergence is controlled
by the Amari-Chentsov third-order tensor $T_{ijk} = \E[\partial_i \log p \cdot \partial_j \log p \cdot \partial_k \log p]$.
Bounds on this tensor for empirical SCX output distributions are needed.
\end{quote}

\vspace{1em}
\noindent\rule{\textwidth}{0.5pt}
\vspace{0.5em}

\noindent {\small \textbf{审查依据：} 本更正基于以下审查文件：
\texttt{docs/reviews/GAUGE\_5VIEWPOINTS\_FINAL.md}（裁决：``夸大''）及
\texttt{docs/reviews/GAUGE\_VIEWPOINTS\_REVIEW.md}（详细分析）。
对观点4的裁决为$\checkmark$数学关系正确但``自然涌现''严重夸大，需降级为``条件等价''。}

\noindent {\small \textbf{Review basis:} This corrigendum is based on:
\texttt{docs/reviews/GAUGE\_5VIEWPOINTS\_FINAL.md} (verdict: ``overstated'') and
\texttt{docs/reviews/GAUGE\_VIEWPOINTS\_REVIEW.md} (detailed analysis).
The verdict on Viewpoint~4 is: $\checkmark$ mathematically correct relationship but
``natural emergence'' is severely overstated — downgrade to ``conditional equivalence''.}

\end{document}
